\documentclass[10pt]{article}
\usepackage[margin=0.65in]{geometry}
\usepackage[T1]{fontenc}
\usepackage[utf8]{inputenc}
\usepackage{lmodern}
\usepackage{enumitem}
\usepackage{graphicx}
\usepackage{eso-pic}
\usepackage{xcolor}
\usepackage[hidelinks]{hyperref}
\setlength{\parindent}{0pt}
\setlength{\parskip}{0.42em}
\newcommand{\sectiontitle}[1]{%
  \vspace{0.35em}
  {\Large\bfseries #1\par}
  \vspace{0.2em}
  \hrule
  \vspace{0.5em}
}
\newcommand{\CodexWatermark}{%
  \AddToShipoutPictureBG{%
    \AtPageCenter{%
      \makebox(0,0){%
        \rotatebox{45}{%
          \textcolor[gray]{0.93}{\fontsize{22}{22}\selectfont\bfseries Winston Huang-Generated with Codex}%
        }%
      }%
    }%
  }%
}
\begin{document}
\CodexWatermark
{\centering
{\LARGE\bfseries AP Microeconomics Curve Challenge\par}
\vspace{0.3em}
{\large Answer Key and Explanations\par}
\vspace{0.8em}}
\sectiontitle{Solutions}
\begin{enumerate}[series=curves, leftmargin=*, itemsep=0.58em, topsep=0.18em, parsep=0pt]
\item \textbf{Answer: C}. Because supply is perfectly inelastic at $Q_f$, a rightward shift in demand raises price but cannot increase quantity. The new equilibrium is at price $P_2$ and quantity $Q_f$.
\item \textbf{Answer: B}. With fixed seating, quantity supplied remains $Q_f$. At the ceiling price $P_c$, quantity demanded is larger than $Q_f$, so a shortage results even though the number of seats sold stays fixed at $Q_f$.
\item \textbf{Answer: A}. Total revenue is price times quantity. Quantity stays fixed at $Q_f$, while price rises from $P_1$ to $P_2$, so revenue increases.
\item \textbf{Answer: D}. A binding ceiling below equilibrium encourages nonprice rationing and can create resale or scalping. Some consumers may ultimately pay above the legal ceiling in secondary transactions.
\end{enumerate}
\begin{enumerate}[resume*=curves, leftmargin=*, itemsep=0.58em, topsep=0.18em, parsep=0pt]
\item \textbf{Answer: E}. At the free-trade world price $P_w$, domestic quantity supplied is $Q_1$ and domestic quantity demanded is $Q_4$. Imports therefore equal $Q_4-Q_1$.
\item \textbf{Answer: B}. At the tariff-inclusive price $P_t$, domestic supply rises to $Q_2$ and domestic demand falls to $Q_3$. Imports shrink to $Q_3-Q_2$.
\item \textbf{Answer: B}. Tariff revenue equals the tariff per unit times the number of imported units after the tariff. That is $(P_t-P_w)(Q_3-Q_2)$.
\item \textbf{Answer: C}. The higher domestic price benefits domestic producers, and the government collects tariff revenue. Domestic consumers are harmed because they pay a higher price and buy less.
\end{enumerate}
\begin{enumerate}[resume*=curves, leftmargin=*, itemsep=0.58em, topsep=0.18em, parsep=0pt]
\item \textbf{Answer: A}. The graph shows the consumer-price increase is larger than the producer-price decrease, so consumers bear the larger share. Tax incidence depends on elasticities, not on the legal assignment of the tax.
\item \textbf{Answer: C}. Tax revenue is the per-unit tax wedge multiplied by the quantity sold after the tax. The vertical wedge is $P_c-P_p$, and the post-tax quantity is $Q_t$.
\item \textbf{Answer: A}. The tax eliminates mutually beneficial trades from $Q_t$ to $Q^{*}$. The deadweight loss is the triangle with height equal to the tax wedge and base equal to the lost quantity.
\item \textbf{Answer: A}. When supply is more elastic, producers are more able to reduce quantity supplied in response to the tax. The less elastic side of the market then bears more of the burden, so consumers pay a larger share.
\end{enumerate}
\begin{enumerate}[resume*=curves, leftmargin=*, itemsep=0.58em, topsep=0.18em, parsep=0pt]
\item \textbf{Answer: A}. A competitive firm chooses the quantity where price equals marginal cost on the upward-sloping portion of MC. At $P_1$, that output is $Q_1$, and the price line lies above ATC there, so economic profit is positive.
\item \textbf{Answer: B}. The shutdown rule compares price with AVC, not ATC. Since $P_2$ lies above AVC and intersects MC at $Q_2$, the firm continues producing $Q_2$ in the short run while taking a loss.
\item \textbf{Answer: C}. A change in fixed cost shifts ATC but does not shift MC or AVC. Because the firm's output rule is $P=MC$ as long as price covers AVC, the output remains $Q_2$.
\item \textbf{Answer: B}. Positive economic profit attracts entry. Entry shifts market supply right and drives the market price downward until firms earn zero economic profit in the long run.
\end{enumerate}
\begin{enumerate}[resume*=curves, leftmargin=*, itemsep=0.58em, topsep=0.18em, parsep=0pt]
\item \textbf{Answer: B}. A single-price monopolist produces where $MR=MC$, which is $Q_m$, and then charges the price consumers are willing to pay on the demand curve, which is $P_m$.
\item \textbf{Answer: B}. Allocative efficiency requires price to equal marginal cost. At the monopoly quantity $Q_m$, the price on the demand curve is above MC, so the monopoly underproduces relative to the efficient level.
\item \textbf{Answer: C}. A perfectly price-discriminating monopolist sells every unit for which willingness to pay is at least marginal cost. That expands output to $Q_c$, where demand intersects MC, the same quantity as in perfect competition.
\item \textbf{Answer: C}. Average-total-cost pricing increases output relative to an unregulated monopoly because price must fall. But since the regulated price still exceeds MC at the chosen quantity, output remains below the allocatively efficient level $Q_c$.
\end{enumerate}
\begin{enumerate}[resume*=curves, leftmargin=*, itemsep=0.58em, topsep=0.18em, parsep=0pt]
\item \textbf{Answer: B}. In long-run monopolistic competition, entry drives economic profit to zero. Graphically, that appears as the demand curve tangent to ATC at the firm's chosen quantity.
\item \textbf{Answer: B}. Excess capacity means the firm produces less than the output that minimizes average total cost. The graph shows the firm's output $Q_m$ is to the left of $Q_e$, the minimum-ATC quantity.
\item \textbf{Answer: C}. Allocative efficiency requires $P=MC$. In monopolistic competition, the downward-sloping demand curve implies price exceeds marginal cost at the profit-maximizing output.
\item \textbf{Answer: A}. Positive short-run profits attract entry of close substitutes. That reduces demand for this firm's product and shifts its demand curve left until economic profit is eliminated.
\end{enumerate}
\begin{enumerate}[resume*=curves, leftmargin=*, itemsep=0.58em, topsep=0.18em, parsep=0pt]
\item \textbf{Answer: B}. A rightward shift in labor demand raises the equilibrium wage and the equilibrium quantity of labor employed in a competitive market.
\item \textbf{Answer: B}. With a binding minimum wage, firms hire only the quantity of labor demanded, which is $Q_d$. Workers supply $Q_s$, so unemployment equals $Q_s-Q_d$.
\item \textbf{Answer: C}. Total earnings for employed workers equal the wage times the number of workers actually employed. Under the binding minimum wage, that is $W_m$ times $Q_d$.
\item \textbf{Answer: B}. If labor demand is relatively inelastic, the wage increase outweighs the reduction in employment, making total earnings more likely to rise.
\end{enumerate}
\begin{enumerate}[resume*=curves, leftmargin=*, itemsep=0.58em, topsep=0.18em, parsep=0pt]
\item \textbf{Answer: B}. A monopsonist hires labor where $MRP=MFC$, which determines $L_m$. It then pays the wage found on the labor supply curve at that employment level, which is $W_m$.
\item \textbf{Answer: C}. In the competitive benchmark, wage and employment are $W_c$ and $L_c$. The monopsonist instead chooses $W_m$ and $L_m$, both of which are lower.
\item \textbf{Answer: C}. A minimum wage set at the competitive wage can flatten the firm's marginal factor cost up to the relevant range. The monopsonist then hires labor where that wage equals MRP, increasing employment to $L_c$.
\item \textbf{Answer: C}. In a monopsony, to attract another worker the employer generally must offer a higher wage, and that higher wage applies to all previously hired workers. Therefore the marginal cost of an extra worker exceeds the wage itself.
\end{enumerate}
\begin{enumerate}[resume*=curves, leftmargin=*, itemsep=0.58em, topsep=0.18em, parsep=0pt]
\item \textbf{Answer: C}. Private firms ignore external costs and produce where demand intersects MPC, which is the market quantity $Q_m$. The socially efficient quantity is lower, at $Q_s$.
\item \textbf{Answer: B}. The corrective tax must internalize the marginal external cost at the socially efficient output. Graphically, that is the vertical gap between $MSC$ and $MPC$ at $Q_s$.
\item \textbf{Answer: C}. Overproduction from $Q_s$ to $Q_m$ creates a triangle of lost welfare. On those excess units, marginal social cost exceeds marginal social benefit.
\item \textbf{Answer: B}. The tax reduces overproduction and moves the market to the socially efficient quantity $Q_s$. Eliminating the deadweight loss increases total surplus.
\end{enumerate}
\begin{enumerate}[resume*=curves, leftmargin=*, itemsep=0.58em, topsep=0.18em, parsep=0pt]
\item \textbf{Answer: C}. A Lorenz curve closer to the 45-degree line indicates a more equal distribution of income. Country B's curve is consistently closer to that line than Country A's.
\item \textbf{Answer: B}. The Gini coefficient rises as the Lorenz curve bows farther away from the line of equality. Since Country A is more bowed out, Country B has the smaller Gini coefficient.
\item \textbf{Answer: B}. Redistribution that makes lower-income households receive a larger share of total income moves the Lorenz curve closer to the line of equality.
\item \textbf{Answer: D}. If the bottom 80\% receive 50\%, the remaining top 20\% must receive the other 50\% of total income.
\end{enumerate}
\end{document}
